\uuid{7cR4}
\titre{QCM — Séries à termes positifs }

\niveau{L1}
\module{Séries numériques}
\chapitre{Séries à termes positifs}
\sousChapitre{Comparaison, équivalents et séries usuelles}
\theme{séries numériques}
\auteur{Quentin Liard}
\datecreate{ 2026-04-10}
\organisation{AMSCC}
\difficulte{3}

\contenu{

\texte{ 
Pour chacune des affirmations suivantes, dire si elle est vraie ou fausse, en justifiant.
}

\begin{enumerate}

\item   \question{La série
$
\sum\limits_{n\geq 1}\ln\left(1+\frac1n\right)
$
converge.}
\indication{}
\reponse{
\textbf{Faux.} On a
$
\ln\left(1+\frac1n\right)=\ln\left(\frac{n+1}{n}\right)=\ln(n+1)-\ln n.
$
La série est donc télescopique. Sa somme partielle d’ordre $N$ vaut
$
\sum\limits_{n=1}^N \ln\left(1+\frac1n\right)=\ln(N+1),
$
qui tend vers $+\infty$. La série diverge.
}




\item   \question{Si $u_n\underset{n\to+\infty}{\sim} v_n$ avec pour tout $n \geq 3$, $u_n\ge 0$, $v_n\ge 0$ et que la série $\sum\limits_{n\geq 0} u_n$ diverge alors la série $\sum\limits_{n\geq 0}   v_n$ diverge.}
\indication{}
\reponse{
\textbf{Vrai.} C’est le théorème de comparaison par équivalence pour les séries à termes positifs.
}


\item   \question{La série
$
\sum\limits_{n\geq 1}\frac{e^{1/n}-1}{n}
$
converge.}
\indication{}
\reponse{
\textbf{Vrai.} Comme
$
e^x-1\sim x \quad \text{quand } x\to 0,
$
on obtient
$
e^{1/n}-1\sim \frac1n.
$
Ainsi
$
\frac{e^{1/n}-1}{n}\sim \frac{1}{n^2}.
$
La série est donc de même nature que une série de Riemann convergente.
}

\item   \question{La série
$
\sum\limits_{n\geq 1}\frac{n}{n^2+1}
$
converge.}
\indication{}
\reponse{
\textbf{Faux.} On a
$
\frac{n}{n^2+1}\sim \frac1n.
$
La série est donc de même nature que la série harmonique, qui diverge.
}


\item   \question{La série
$
\sum\limits_{n\geq 2}\left(\frac{2}{3}\right)^n
$
converge et a pour somme $\dfrac{4}{3}$.}
\indication{}
\reponse{
\textbf{Vrai.} Il s'agit d'une série géométrique de raison $r=\frac{2}{3}$, qui converge car $|r|<1$. En utilisant la formule de la somme d'une série géométrique convergente à partir de l'indice $n=0$~:
$
\sum_{n=0}^{+\infty}\left(\frac{2}{3}\right)^n = \frac{1}{1-\frac{2}{3}}=3,
$
on en déduit
$
\sum_{n=2}^{+\infty}\left(\frac{2}{3}\right)^n
= 3 - \left(\frac{2}{3}\right)^0 - \left(\frac{2}{3}\right)^1
= 3 - 1 - \frac{2}{3} = \frac{4}{3}.
$
}

\end{enumerate}

}